\section{Predictions and a research agenda}\label{sec:agenda}
The argument yields predictions that could fail:
\begin{enumerate}
\item \textbf{P1 (moderation).} Baseline executive function, measured before treatment, predicts the acquisition of restructuring skill more than the acquisition of exposure-based skills.
\item \textbf{P2 (stress-dependence).} Cognitive reappraisal efficacy falls with stress intensity and with lower executive capacity, so the gap between ``calm'' and ``in-the-moment'' use is larger in low-executive-function samples.
\item \textbf{P3 (component).} In adult ADHD with anxiety, a behaviour-and-exposure-only protocol is not inferior to the same protocol plus restructuring.
\item \textbf{P4 (ADHD).} ADHD status predicts less benefit from the cognitive module specifically, not from the skills modules.
\end{enumerate}
\textbf{Evidence that would refute the argument:} an adequately powered study finding no moderation by executive function; component trials showing that restructuring adds substantial benefit; or meta-regression showing that executive-function-heavy samples respond as well. A secondary analysis of existing trial data, or an individual-participant meta-analysis, is the cheapest next step (see \texttt{experiments/}).

\section{Conclusion}
Restructuring-heavy CBT is effective on average. What is not established is that the cognitive component is the active ingredient, or that it works for people whose executive function is low. The best-supported claim is narrow: the argue-with-your-thoughts component is the part least supported mechanistically and most exposed to executive limits, while exposure, skills and practice look like the parts that carry the benefit.
